% Generated draft: requires mathematical and source review.
\documentclass[11pt]{article}
\usepackage[margin=25mm]{geometry}
\usepackage{fontspec}
\setmainfont{Latin Modern Roman}
\usepackage{amsmath,amssymb,mathtools,unicode-math}
\setmathfont{Latin Modern Math}
\usepackage{enumitem,needspace,xurl,hyperref}
\setlength{\parindent}{0pt}
\setlength{\parskip}{5pt}
\setlength{\emergencystretch}{3em}
\newcommand{\N}{\mathbb{N}}
\newcommand{\Z}{\mathbb{Z}}
\newcommand{\Q}{\mathbb{Q}}
\newcommand{\R}{\mathbb{R}}
\newcommand{\C}{\mathbb{C}}
\newcommand{\F}{\mathbb{F}}
\newcommand{\A}{\mathbb{A}}
\newcommand{\PP}{\mathbb{P}}
\newcommand{\E}{\mathbb{E}}
\newcommand{\eps}{\varepsilon}
\newcommand{\dd}{\,\mathrm{d}}
\begin{document}
\section*{UnsolvedMath ID: 700549}
\textbf{Source problem code:} catalogue entry
\subsection{Definitions and mathematical statement}
A group \(G\) is \emph{locally finite} if \(\langle S\rangle\) is finite for every finite subset \(S\subseteq G\). A subgroup \(A\leq G\) is \emph{periodic} if every \(a\in A\) has finite order, and is \emph{abelian} if its elements commute with one another.

For subgroups \(A,B\leq G\), write \(AB=\{ab:a\in A,\ b\in B\}\). There are two potentially different meanings of subgroups that ``commute'': setwise permutability, \(AB=BA\), and elementwise commutation, \(ab=ba\) for all \(a\in A\), \(b\in B\). The supplied wording does not explicitly distinguish them.

Under the finite-product reading and the interpretation of commuting as setwise permutability, the question has the following precise candidate formulation: for every integer \(m\geq1\), every group \(G\), and all periodic abelian subgroups \(A_1,\ldots,A_m\leq G\) satisfying
\[
G=A_1A_2\cdots A_m,
\qquad
A_iA_j=A_jA_i\quad(1\leq i,j\leq m),
\]
must \(G\) be locally finite? Here \(A_1\cdots A_m=\{a_1\cdots a_m:a_i\in A_i\}\).

If elementwise commutation is intended, the pairwise condition instead reads
\[
\forall i,j\in\{1,\ldots,m\}\ \forall a\in A_i\ \forall b\in A_j,
\qquad ab=ba.
\]
These interpretations are recorded separately; the supplied text does not establish which is intended. It also does not explicitly specify whether the number of factors must be finite. Thus the displayed finite-product formulation is conditional on that reading, and an intended infinite-product formulation would require an additional definition of the product.

The source states that the answer is affirmative for two subgroups. This assertion is retained as supplied background, with the same unresolved convention concerning commuting subgroups.
\subsection{Short English statement}
D. I. Zaitsev asks whether a group expressed as a product of pairwise commuting periodic abelian subgroups must be locally finite. The source reports an affirmative answer for two factors. The meaning of commuting subgroups and the permitted number of factors require clarification.
\subsection{Sources}
\begin{itemize}
\item Supplied local UnsolvedMath registry record: permanent ID \(700549\), exact source code \texttt{catalogue entry}, title \emph{Kourovka Notebook Problem 9.14}; upstream revision \texttt{d371ee917ac9647fbb06cd6fc529deec6f3dda82}.
\item \emph{The Kourovka Notebook (Unsolved Problems in Group Theory)}, problem 9.14, attributed to D. I. Zaitsev in the 9th issue (1984). The supplied record also identifies the 21st issue edition (2026), \texttt{arXiv:1401.0300v48}. The collection description names E. I. Khukhro and V. D. Mazurov as editors. These bibliographic details are supplied attribution and have not been independently checked.
\item The registry labels the problem \texttt{open} and states that the two-subgroup case has an affirmative answer. These are supplied status and background assertions, not independently verified mathematical or literature assessments. No separate citation for the two-subgroup assertion is supplied.
\end{itemize}
\end{document}
